
\appendix
\section{Diagram Inventory and Viewpoint Guide}
Each figure is editable TikZ source within the corresponding chapter fragment. The diagrams are conceptual and intentionally avoid vendor-specific deployment assumptions.

\begin{longtable}{@{}p{17mm}p{49mm}p{66mm}@{}}\toprule
\textbf{Figure} & \textbf{Architectural viewpoint} & \textbf{Reviewer question}\\\midrule\endhead
\ref{fig:taxonomy} & Responsibility boundaries & Which layer has execution authority?\\
\ref{fig:chaining} & Sequential fixed workflow & Where are intermediate outputs validated?\\
\ref{fig:parallel} & Static fan-out and fan-in & What happens to missing or conflicting results?\\
\ref{fig:routing} & Conditional path selection & How are ambiguous/high-risk inputs routed?\\
\ref{fig:react} & Dynamic tool loop & What prevents unbounded or forbidden actions?\\
\ref{fig:plan} & Plan-and-execute & How are stale plans invalidated?\\
\ref{fig:supervisor} & Dynamic hierarchical delegation & How are worker scopes and synthesis governed?\\
\ref{fig:collab} & Peer-agent collaboration & Which service checks peer identities and permissions?\\
\ref{fig:evaluator} & Feedback loop & What constitutes independent acceptance evidence?\\
\ref{fig:approval} & Human authority & Is approval bound to the exact current action?\\
\ref{fig:interoperability} & MCP/A2A boundary & What does the protocol provide versus the harness?\\
\ref{fig:appsec} & Composed application reference & Which paths remain read-only and which require review?\\\bottomrule
\end{longtable}

\section{Source Register, Verification, and Limitations}
The sources below are linked to the first-party technical descriptions, foundational research, and governance documents reviewed on 23 November 2026. Comparisons, diagrams, formulas, Python code, and measurable test criteria in this document are original engineering synthesis, not text reproduced from source materials. Standards and protocol revisions are subject to change.

\begin{thebibliography}{99}\small
\bibitem{anthropic} Anthropic, \emph{Building Effective Agents}, 19 December 2024. Defines prompt chaining, routing, parallelization, orchestrator--workers, evaluator--optimizer, and autonomous agents.
\url{https://www.anthropic.com/research/building-effective-agents}

\bibitem{anthropic-guide} Anthropic, \emph{Common Architecture Patterns}, architecture implementation guide.
\url{https://claude.com/resources/guides/building-effective-ai-agents/common-architecture-patterns}

\bibitem{react} Yao et al., \emph{ReAct: Synergizing Reasoning and Acting in Language Models}, ICLR 2023, arXiv:2210.03629.
\url{https://arxiv.org/abs/2210.03629}

\bibitem{databricks} Databricks, Molly Nagamuthu, \emph{Four AI Architecture Patterns for Regulated Industries}, 19 November 2026.
\url{https://www.databricks.com/blog/ai-architecture-patterns-regulated-industries}

\bibitem{mcp} Model Context Protocol, \emph{Architecture}, 2026-07-28 specification snapshot.
\url{https://modelcontextprotocol.io/specification/2026-07-28/architecture}

\bibitem{a2a} A2A Project, \emph{Agent2Agent Protocol Specification}, release 1.0.0, as published at verification.
\url{https://a2a-protocol.org/v1.0.0/specification/}

\bibitem{owasp} OWASP Cheat Sheet Series, \emph{AI Agent Security Cheat Sheet}.
\url{https://cheatsheetseries.owasp.org/cheatsheets/AI_Agent_Security_Cheat_Sheet.html}

\bibitem{nist} NIST, \emph{AI Risk Management Framework (AI RMF 1.0)}.
\url{https://www.nist.gov/itl/ai-risk-management-framework}

\bibitem{nistgenai} NIST, \emph{Artificial Intelligence Risk Management Framework: Generative Artificial Intelligence Profile}, NIST AI 600-1, 2024.
\url{https://doi.org/10.6028/NIST.AI.600-1}
\end{thebibliography}

\section{Repository Build and Maintenance}
From the root of the latex-docs repository, run:
\begin{verbatim}
make preflight
make build-category-data-systems
\end{verbatim}
The source of record is the main LaTeX root, which imports seven version-controlled chapter fragments. The twelve comparative diagrams are embedded as editable TikZ, so the guide does not require an external image asset. Generated PDFs are build products, not the canonical source. Re-run source verification whenever protocol versions or high-impact policy recommendations change.
